% ════════════════════════════════════════════════════════════════
\section{Wykład 2}
\subsection{Notacja asymptotyczna --- przypomnienie}
% ════════════════════════════════════════════════════════════════
\begin{flalign*}
	& f,g:\N\rightarrow\Rp & \\
	& g(n) = \BO(f(n)) \Leftrightarrow \exists_{c\in\Rp}\exists_{n_0\in\Np}\forall_{n > n_0} g(n) \leq c\cdot f(n) & \\
	& g(n) = \BM(f(n)) \Leftrightarrow \exists_{c\in\Rp}\exists_{n_0\in\Np}\forall_{n > n_0} c\cdot f(n) \leq g(n) & \\
	& g(n) = \BT(f(n)) \Leftrightarrow \exists_{c_1, c_2\in\Rp}\exists_{n_0\in\Np}\forall_{n > n_0} c_1\cdot f(n)\leq g(n) \leq c_2\cdot f(n) & \\
	& g(n) = \BO(f(n)) \Leftrightarrow f(n) = \BM(g(n)) & \\
	& g(n) = \BT(f(n)) \Leftrightarrow \big( g(n) = \BO(f(n)) \wedge g(n) = \BM(f(n)) \big) &
\end{flalign*}

\begin{flalign*}
	& n^a = \BO(n^b),\ 0 < a \leq b && \log n = \BO(n^\epsilon),\ \epsilon > 0 & && \\ 
	& a^n = \BO(b^n),\ 1 < a \leq b && \log_a n = \BT(\log_b n),\ a, b > 1 & && \\ 
	& n^a = \BO(b^n),\ 0 < a,\ b > 1 && &
\end{flalign*}

\begin{flalign*}
	T(n) = 5n^3 + 2n^2 + 100n\log n \underset{n \to \infty}{\approx} 5n^3 = \BT(n^3) && && 
\end{flalign*}

\paragraph{Procedura rekurencyjna}
\begin{flalign*}
	n! = \begin{cases}
		1 & n = 0 \\[0.5ex]
		n \cdot (n - 1)! & n \geq 1 
	\end{cases} && && 
\end{flalign*}
Oczywiste jest, że $\displaystyle n! = \prod_{i = 1}^{n}i$

\paragraph{Ciąg Fibonacciego}
	$F_n$ - ciąg Fibonacciego.
\begin{flalign*}
	F_n = \begin{cases}
		0 & n = 0 \\[0.5ex]
		1 & n = 1 \\[0.5ex]
		F_{n - 1} + F_{n - 2} & n \geq 2
	\end{cases} && && 
\end{flalign*}

\begin{alg}{Rekurencyjna konstrukcja ciągu Fibonacciego}{alg:fib:recursive}
	\FunctionNN{Fibb}{$n$}
	\If{$n = 0$} \Return 0 
	\EndIf
	\If{$n = 1$} \Return 1
	\EndIf
	\State \Return \Call{Fibb}{$n - 1$} $+$ \Call{Fibb}{$n - 2$}
	\EndFunction
\end{alg}

Wywołanie rekurencyjne dla (typowo łatwiejszych i mniejszych) instancji tego samego problemu.

\subsection{Strategia typu \textit{Dziel i Zwyciężaj} (\textit{Divide and Conquer})}
Algorytm typu \textit{dziel i zwyciężaj} składa się z trzech części
\begin{enumerate}
	\item \textit{Dziel} -- problem jest dzielony na mniejsze podproblemy
	\item \textit{Zwyciężaj} -- podproblemy są rozwiązywane rekurencyjnie
	\item \textit{Połącz} -- rozwiązanie całego problemu jest konstruowane z rozwiązań podproblemów
\end{enumerate}

Zakładamy, że (zazwyczaj) w fazie \textit{dziel} problem jest dzielony na co najmniej dwa podproblemy i podproblemy są rozłączne.

\begin{example}{Mnożenie liczb $n$-cyfrowych} 
	$x,y$ -- liczby naturalne $n$-cyfrowe.
	
	Aby policzyć $x\cdot y$ można użyć mnożenia pisemnego (jak na kartce). Tym sposobem musimy wykonać $\BO(n^2)$ mnożeń cyfr. Spróbujemy znaleźć lepszą metodę. 
	Dzielimy $x$ i $y$ na ,,dwie połowy''. Zakładamy, dla uproszczenia, że $n$ jest parzyste.
	\begin{flalign*}
		x = x_1x_2\cdots x_n = \underbracket{x_1x_2\cdots x_{n/2}}_{x_L}\underbracket{x_{n/2 + 1}\cdots x_n}_{x_R} = 10^{n/2}x_L + x_R && && 
	\end{flalign*}
	podobnie dla $y$: 
	\begin{flalign*}
		y = \cdots = 10^{n/2}y_L + y_R && &&
	\end{flalign*}
	
	Wymnażamy:
	\begin{flalign*}
		x\cdot y = (10^{n/2}x_L + x_R) \cdot (10^{n/2}y_L + y_R) = 10^nx_Ly_L + 10^{n/2}(x_Ly_R + x_Ry_L) + x_Ry_R && && 
	\end{flalign*}
	
	Złożoność czasowa algorytmu -- ozn. $T(n)$. 
	\begin{itemize}[nosep]
		\item Podział $x, y$ na $x_L,x_R,y_L,y_R$ zajmie $\BO(n)$ czasu.
		\item Dodanie liczb $n$-cyfrowych zajmie $\BO(n)$ czasu.
		\item Mnożenie przez $10^n$ jest równoważne dopisaniu $n$ zer więc też $\BO(n)$.
	\end{itemize}
	\begin{flalign*}
		T(n) = 4\cdot T(n/2) + \BT(n) && && 
	\end{flalign*}
	
	Można lepiej, bo możemy zauważyć, że $x_Ly_R + x_Ry_L = (x_L - x_R)\cdot (y_R - y_L) + x_L y_L + x_R y_R$, 
	więc potrzebujemy policzyć tylko 3 iloczyny: $x_L y_L,\ x_R y_R,\ (x_L - x_R)\cdot (y_R - y_L)$. 
	To daje $T(n) = 3\cdot T\left(\frac{n}{2}\right) + \BT(n)$
\end{example}

\begin{theorem}{CLRS - Uniwersalne Twierdzenie o Rekurencji} \label{theorem:utor}
	Niech $a \geq 1,\ b > 1$ będą stałymi, $f(n)$ funkcją i $T(n) = a\cdot T\left(\frac{n}{b}\right) + f(n)$ (gdzie $\frac{n}{b}$ znaczy $\left\lfloor \frac{n}{b} \right\rfloor$ lub $\left\lceil \frac{n}{b}\right\rceil$).
	Wtedy
	\begin{flalign*}
		T(n) = \begin{cases}
			\BT(n^{\log_b a}) & f(n) = \BO(n^{\log_b a - \epsilon}) \text{ dla pewnego } \epsilon > 0 \\[1ex]
			\BT(n^{\log_b a} \log n) & f(n) = \BT(n^{\log_b a}) \\[1ex]
			\BT(f(n)) & \substack{
				f(n) = \BM(n^{\log_b a + \epsilon}) \text{ dla pewnego } \epsilon > 0 \\ 
				a\cdot f(n/b) \leq c\cdot f(n) \text{ dla pewnej stałej } c < 1 \text{ dla dużych } n
			}
		\end{cases} && && 
	\end{flalign*}
\end{theorem}

\begin{lemma}
	Niech $a\geq 1,\ b>1$ będą stałymi i niech $f(n)$ będzie nieujemną funkcją zdefiniowaną dla potęg $b$. Jeśli
	\begin{flalign*}
		T(n) = \begin{cases}
			\BT(1) & n = 1 \\[1ex]
			a\cdot T\left(\frac{n}{b}\right) + f(n) & n = b^i,\ i \geq 1
		\end{cases} && && 
	\end{flalign*}
	to 
	\begin{flalign*}
		T(n) = \BT\left(n^{\log_b n}\right) + \sum_{j = 0}^{\log_b n - 1} a^j f\left(\frac{n}{b^j}\right) && && 
	\end{flalign*}
\end{lemma}

\begin{proof}
	\begin{flalign} \label{eq:utor:star}
		a^{\log_b n} = b^{\log_b n \cdot \log_b a} = b^{\log_b a \cdot \log_b n} = \left(b^{\log_b a}\right)^{\log_b n} = n^{\log_b a} && && 
	\end{flalign}
	
	\begin{align*}
		T(n) &= f(n) + a\cdot T\left(\frac{n}{b}\right) 
		= f(n) + a\left(f\left(\frac{n}{b}\right) + aT\left(\frac{n}{b^2}\right)\right) \\[1.5ex]
		&= f(n) + af\left(\frac{n}{b}\right) + a^2 T\left(\frac{n}{b^2}\right) 
		= f(n) + af\left(\frac{n}{b}\right) + a^2\left(f\left(\frac{n}{b^2}\right) + aT\left(\frac{n}{b^3}\right)\right) \\[1.5ex]
		&= f(n) + af\left(\frac{n}{b}\right) + a^2 f\left(\frac{n}{b^2}\right) + \cdots + a^{\log_b n - 1} f\left(\frac{n}{b^{\log_b n}}\right) + a^{\log_b n} T(1) \\[1.5ex]
		&\overset{\ref{eq:utor:star}}{=} \cdots + a^{\log_b n} T(1) 
		= \cdots + n^{\log_b a} T(1) 
		= \BT\left(n^{\log_b a}\right) + \sum_{j = 0}^{\log_b n - 1} a^j f\left(\frac{n}{b^j}\right)
	\end{align*}
\end{proof}

Przeprowadzimy teraz dowód uniwersalnego twierdzenia o rekurencji (\ref{theorem:utor}). Ograniczymy się do przypadku, gdy $n$ jest potęgą $b$. Pozostały przypadek jest bardziej techniczny.

\begin{proof}
	Rozważmy wszystkie przypadki:
	\begin{enumerate}
		\item $f(n) = \BO\left(n^{\log_b a - \epsilon}\right)$. Wtedy
		\begin{flalign*}
			f\left(\frac{n}{b^j}\right) &= \BO\left(\left(\frac{n}{b^j}\right)^{\log_b a - \epsilon}\right) && && \\[1.5ex]
			%
			\sum_{j = 0}^{\log_b n - 1} a^j f\left(\frac{n}{b^j}\right) 
			&= \sum_{j = 0}^{\log_b n - 1} a^j \BO\left(\left(\frac{n}{b^j}\right)^{\log_b a - \epsilon}\right) 
			= \BO\left(\underbracket{\sum_{j = 0}^{\log_b n - 1} a^j \left(\frac{n}{b^j}\right)^{\log_b a - \epsilon}}_{\blacktriangle}\right) && && \\[1.5ex]
			%
			\blacktriangle &= \sum_{j = 0}^{\log_b n - 1} a^j \frac{n^{\log_b a - \epsilon}}{b^{j(\log_b a - \epsilon)}} 
			= n^{\log_b a - \epsilon} \sum_{j = 0}^{\log_b n - 1} a^j \frac{b^{j\epsilon}}{b^{j\log_b a}} && && \\[1.5ex]
			&= n^{\log_b a - \epsilon} \sum_{j = 0}^{\log_b n - 1} \left(\frac{a\cdot b^{\epsilon}}{\underbracket{b^{\log_b a}}_{= a}}\right)^j 
			= n^{\log_b a - \epsilon} \sum_{j = 0}^{\log_b n - 1} b^{\epsilon j} \quad \text{(wzór na ciąg geometryczny)} && && \\[1.5ex]
			&= n^{\log_b a - \epsilon} \frac{b^{\epsilon \log_b n} - 1}{b^\epsilon - 1} 
			= n^{\log_b a - \epsilon} \frac{n^\epsilon - 1}{\underbracket{b^\epsilon - 1}_{\text{stała}}} && && \\[1.5ex]
			&= n^{\log_b a - \epsilon} \BO(n^\epsilon - 1) 
			= \BO\left(n^{\log_b a - \epsilon}\cdot n^\epsilon - 1\right) 
			= \BO\left(n^{\log_b a}\right)  && && \\[1.5ex]
			%
			\text{Z lematu: } T(n) &= \sum_{j = 0}^{\log_b n - 1} a^j T\left(\frac{n}{b^j}\right) + \BT\left(n^{\log_b a}\right) 
			= \BO\left(n^{\log_b a}\right) + \BT\left(n^{\log_b a}\right) 
			= \BT\left(n^{\log_b a}\right) && &&
		\end{flalign*}
		
		\item $f(n) = \BT\left(n^{\log_b a}\right)$. Wtedy:
		\begin{flalign*}
			f\left(\frac{n}{b^j}\right) &= \BT\left(\left(\frac{n}{b^j}\right)^{\log_b a}\right) && && \\[1.5ex]
			%
			\sum_{j = 0}^{\log_b n - 1} a^j f\left(\frac{n}{b^j}\right) 
			&= \sum_{j = 0}^{\log_b n - 1} a^j \BT\left(\left(\frac{n}{b^j}\right)^{\log_b a}\right) 
			= \BT\left(\sum_{j = 0}^{\log_b n - 1} a^j \left(\frac{n}{b^j}\right)^{\log_b a}\right) && && \\[1.5ex]
			&= \BT\left(n^{\log_b a} \sum_{j = 0}^{\log_b n - 1} a^j \frac{1}{b^{j\log_b a}}\right) 
			= \BT\left(n^{\log_b a} \sum_{j = 0}^{\log_b n - 1} a^j \frac{1}{\underbracket{b^{\log_b a^j}}_{=a^j}}\right) && && \\[1.5ex]
			&= \BT\left(n^{\log_b a} \sum_{j = 0}^{\log_b n - 1} 1 \right) 
			= \BT\left(n^{\log_b a} \log_b n\right) && && \\[1.5ex]
			%
			\text{Z lematu: } T(n) &= \sum_{j = 0}^{\log_b n - 1} a^j f\left(\frac{n}{b^j}\right) + \BT\left(n^{\log_b a}\right) 
			= \BO\left(n^{\log_b a}\log n\right) + \BT\left(n^{\log_b a}\right) 
			= \BT\left(n^{\log_b a}\log n\right) && &&
		\end{flalign*}
		
		\item $f(n) = \BM\left(n^{\log_b a + \epsilon}\right)$ oraz $a \cdot f\left(\frac{n}{b}\right) \leq c\cdot f(n)$ dla jakiegoś $c < 1$ i odpowiednio dużych $n$.
		
		Z założenia $f\left(\frac{n}{b}\right) \leq \frac{c}{a}f(n)$, więc $f\left(\frac{n}{b^j}\right) \leq \left(\frac{c}{a}\right)^j f(n) = \frac{c^j}{a^j}f(n)$.
		
		\begin{flalign*}
			\sum_{j = 0}^{\log_b n - 1} a^j f\left(\frac{n}{b^j}\right) 
			&\leq \sum_{j = 0}^{\log_b n - 1} \cancel{a^j} \frac{c^j}{\cancel{a^j}} f(n) + \BO(1) 
			\leq f(n) \sum_{j = 0}^{\log_b n - 1} c^j + \BO(1) && && \\[1.5ex]
			&\leq f(n) \sum_{j = 0}^{\infty} c^j + \BO(1) 
			\leq f(n) \frac{1}{1 - c} + \BO(1) 
			= \BO(f(n)) && && \\[1.5ex]
			%
			\text{Z drugiej strony: } \sum_{j = 0}^{\log_b n - 1} a^j f\left(\frac{n}{b^j}\right) 
			&= a^0 f\left(\frac{n}{b^0}\right) + \underbracket{\sum_{j = 1}^{\log_b n - 1} a^j f\left(\frac{n}{b^j}\right)}_{\geq 0} 
			= \BM(f(n)) && && \\[1.5ex]
			%
			\text{Z lematu: } T(n) &= \sum_{j = 0}^{\log_b n - 1} a^j f\left(\frac{n}{b^j}\right) + \BT(n^{\log_b a}) 
			= \BT(f(n)) + \BT(n^{\log_b a}) = \BT(f(n)) && &&
		\end{flalign*}
	\end{enumerate}
\end{proof}

\begin{conclusion}
	Niech $a\geq 1,\ b>1$ będą stałymi, a $f(n)$ funkcją taką, że $f(n) = \BT(n^k)$.
	Niech $T(n) = a\cdot T\left(\frac{n}{b}\right) + f(n)$ (gdzie $\frac{n}{b}$ znaczy $\left\lfloor \frac{n}{b} \right\rfloor$ lub $\left\lceil \frac{n}{b}\right\rceil$).
	Wtedy
	\begin{flalign*}
		& T(n) = \begin{cases}
			\BT(n^{\log_b a}) & a > b^k \Leftrightarrow \log_b a > k \\[1ex]
			\BT(n^{\log_b a}\log n) & a = b^k \Leftrightarrow \log_b a = k \\[1ex]
			\BT(n^k) & a < b^k \Leftrightarrow \log_b a < k
		\end{cases} && &&
	\end{flalign*}
\end{conclusion}

Wróćmy teraz do naszego problemu mnożenia liczb $n$-cyfrowych. 
Drugi algorytm ma złożoność:
\begin{flalign*}
	T(n) &= 4\cdot T\left(\frac{n}{2}\right) + \BT(n^1) && a = 4,\ b = 2,\ k = 1 && \\[1.5ex]
	4 &> 2^1 \Rightarrow T(n) = \BT\left(n^{\log_2 4}\right) = \BT\left(n^2\right) && &&
\end{flalign*}

Trzeci algorytm ma złożoność:
\begin{flalign*}
	T(n) &= 3\cdot T\left(\frac{n}{2}\right) + \BT(n^1) && a = 3,\ b = 2,\ k = 1 && \\[1.5ex]
	3 &> 2^1 \Rightarrow T(n) = \BT\left(n^{\log_2 3}\right) = \BT\left(n^{1.59}\right) && &&
\end{flalign*}